\section{Analytic lemmas underlying the linear certificates}
\label{supp:proof-lemmas}
This section spells out the analytical implications tested by the
certificates. It concerns three spatial dimensions and radial angular
sectors $\ell=0,1$ only. A radial reduction of the three-dimensional
problem is not a one-dimensional field theory. The normalization,
origin powers and exterior kernels below depend on dimension and
angular momentum; they cannot be transferred unchanged to other cases.
All inequalities are evaluated over the whole dyadic parameter box.

\subsection{Profile tail: contraction and matching errors}
\label{supp:profile-tail}
Let $1/\sqrt2<\omega<1$, $\kappa_0=\sqrt{1-\omega^2}$, $L,c,\bar f>0$,
and define
\begin{align}
 K_P&=\frac{(2+\tfrac32\bar f^2)e^{-2\kappa_0L}}{4\kappa_0^2L^2},
 &\eta&=2K_Pc^3,&m&=c+\eta,\nonumber\\
 \Lambda&=\frac{(6+\tfrac{15}{2}\bar f^2)m^2e^{-2\kappa_0L}}
                   {4\kappa_0^2L^2}.&&&
 \label{eq:proof-profile-constants}
\end{align}
Assume $m e^{-\kappa_0L}/L\leq\bar f$, $K_Pm^3\leq\eta$,
$\Lambda<1$ and $c>\eta$. Then the profile Volterra equation has a
unique solution in $\|\widetilde u-c\|_\infty\leq\eta$,
where $f=e^{-\kappa_0r}\widetilde u/r$. It is positive and solves the
profile equation. With $\eta_0=K_Pm^3$, its matching errors satisfy
\begin{equation}
 |Le^{\kappa_0L}f(L)-c|\leq\eta_0,\qquad
 |e^{\kappa_0L}(Lf'(L)+f(L))+\kappa_0c|\leq2\kappa_0\eta_0.
 \label{supp:profile-errors}
\end{equation}

\noindent\emph{Proof.}
On the indicated complete ball in $C_b([L,\infty))$, $|\widetilde u|\leq m$
and $|f|\leq\bar f$. The kernel is bounded by $1/(2\kappa_0)$ and
\begin{equation}
 \int_r^\infty e^{-2\kappa_0s}s^{-2}\,\dd s
 \leq\frac{e^{-2\kappa_0r}}{2\kappa_0r^2}.
\end{equation}
Thus the integral map differs from $c$ by at most $K_Pm^3$. The
nonlinear integrand as a function of $\widetilde u$ is
$-2e^{-2\kappa_0s}s^{-2}\widetilde u^3+
\tfrac32e^{-4\kappa_0s}s^{-4}\widetilde u^5$; its derivative is bounded
in modulus by $(6+\tfrac{15}{2}\bar f^2)m^2e^{-2\kappa_0s}s^{-2}$.
The map is therefore a contraction of constant $\Lambda$.
Banach's theorem supplies its fixed point and $\widetilde u\geq c-\eta>0$.
Writing $u=e^{-\kappa_0r}\widetilde u$ gives
\begin{equation}
 u(r)=ce^{-\kappa_0r}+\int_r^\infty
 \frac{\sinh(\kappa_0(s-r))}{\kappa_0}q_f(s)u(s)\,\dd s,
 \qquad q_f=-2f^2+\tfrac32f^4.
\end{equation}
Differentiation proves $u''-\kappa_0^2u=q_fu$ and the profile equation.
The first bound in \eqref{supp:profile-errors} follows from the
self-map estimate; for the second use
\begin{equation}
 e^{\kappa_0r}u'(r)+\kappa_0c
 =-\int_r^\infty\frac{1+e^{-2\kappa_0(s-r)}}{2}
                   q_f(s)\widetilde u(s)\,\dd s.
\end{equation}
The same integral majorant gives $2\kappa_0\eta_0$ at $L$.
The integral tails vanish at infinity, proving the asserted normalization.
For a fixed profile-row scale $s_1$ and $E_0=e^{-\kappa_0L}$,
the resulting matching-map corrections satisfy
\begin{equation}
 |E_1|\leq |s_1|E_0\eta_0/L,\qquad
 |E_2|\leq |s_1|E_0(2\kappa_0\eta_0/L+\eta_0/L^2).
\end{equation}

\subsection{Decaying mode: the two angular sectors}
\label{supp:jost-tail}
Put $dV=-4f^2+\tfrac92 f^4$, $C=-2f^2+3f^4$ and
$q_J=|dV|+|C|\leq6f^2+\tfrac{15}{2}f^4$. A common integral bound is
\begin{equation}
 Q=\frac{(6+\tfrac{15}{2}\bar f^2)m^2e^{-2\kappa_0L}}
          {2\kappa_0L^2}\geq\int_L^\infty q_J(s)\,\dd s.
\end{equation}
Assume $k=\sqrt{(\omega+\rho)^2-1}>0$ and
$\kappa_c=\sqrt{1-(\omega-\rho)^2}>0$. Normalize
$(\widetilde A,\widetilde B)=e^{\kappa_cr}(A,B)$.
For $\ell=0$ their Volterra equation has free term $(0,1)$ and kernels
\begin{equation}
 K_A=e^{-\kappa_c(s-r)}\frac{\sin(k(s-r))}{k},\qquad
 K_B=\frac{1-e^{-2\kappa_c(s-r)}}{2\kappa_c}.
 \label{eq:proof-l0-kernels}
\end{equation}
The first kernel multiplies $dV\widetilde A+C\widetilde B$ and the
second multiplies $dV\widetilde B+C\widetilde A$.
With $\mu=\max(1/k,1/(2\kappa_c))$ and $e_0=e^{\mu Q}-1$, one obtains
\begin{align}
 |\widetilde A|&\leq\frac{e_0}{\mu k},&
 |\widetilde B-1|&\leq\frac{e_0}{2\mu\kappa_c},\nonumber\\
 |e^{\kappa_cr}A'|&\leq
  \left(\sqrt{1+\kappa_c^2/k^2}+\kappa_c/k\right)\frac{e_0}{\mu},&
 |e^{\kappa_cr}B'+\kappa_c|&\leq\frac{3e_0}{2\mu}.
 \label{supp:jost-l0-errors}
\end{align}
These bounds hold for all $r\geq L$.

For $\ell=1$ define the free Riccati functions
\begin{align}
 j(x)&=\sin(x)/x-\cos(x),& y(x)&=-\cos(x)/x-\sin(x),\nonumber\\
 d(u)&=e^{-u}(1+1/u),&g(u)&=e^u(1-1/u),\nonumber\\
 G_A(r,s)&=\frac{j(kr)y(ks)-y(kr)j(ks)}{k},&
 G_B(r,s)&=\frac{d(\kappa_cr)g(\kappa_cs)-g(\kappa_cr)d(\kappa_cs)}{2\kappa_c}.
 \label{eq:proof-l1-kernels}
\end{align}
Here $K_{A,B}=e^{-\kappa_c(s-r)}G_{A,B}$ and the free term is
$(0,1+1/(\kappa_cr))$. For $kL,\kappa_cL\geq1$ put
\begin{align}
 a_1&=\frac{1+(kL)^{-2}}{k},&a_2&=\sqrt{1+(kL)^{-2}},\nonumber\\
 b_0&=1+(\kappa_cL)^{-1},&b_1&=\frac{b_0}{2\kappa_c},&
 b_2&=1+(\kappa_cL)^{-1}+\tfrac12(\kappa_cL)^{-2}.
 \label{eq:proof-l1-majorants}
\end{align}
Then $|K_A|\leq a_1$, $|K_B|\leq b_1$,
$|e^{-\kappa_c(s-r)}\partial_rG_A|\leq a_2$ and
$|(\partial_r-\kappa_c)K_B|\leq b_2$.
With $\mu=\max(a_1,b_1)$ and $e_1=b_0(e^{\mu Q}-1)/\mu$ the four
errors in normalized physical Cauchy data are bounded by
$(a_1,b_1,a_2,b_2)e_1$ about
\begin{equation}
 J_0=\left(0,1+\frac1{\kappa_cL},0,
              -\kappa_c-\frac1L-\frac1{\kappa_cL^2}\right).
\end{equation}
In particular $e^{\kappa_cr}(|A|+|B|)\leq b_0+(a_1+b_1)e_1$.

\noindent\emph{Proof.}
For $\ell=0$, the kernels have bounds $1/k$ and $1/(2\kappa_c)$.
The $n$th ordered Volterra iterate is bounded by $(\mu Q)^n/n!$.
Thus the series converges uniformly and
$\psi=\max(|\widetilde A|,|\widetilde B|)\leq e^{\mu Q(r)}$,
where $Q(r)=\int_r^\infty q_J$. Since
$\int_r^\infty q_J(s)e^{\mu Q(s)}\,\dd s=(e^{\mu Q(r)}-1)/\mu$,
the value bounds follow. Differentiating the kernels gives
$|\partial_rK_A|\leq\sqrt{1+\kappa_c^2/k^2}$ and
$|\partial_rK_B|\leq1$; subtracting $\kappa_c$ times the normalized
values gives the displayed conservative physical derivative bounds.
Twice differentiating the integral equations proves the radial ODE.

For $\ell=1$, direct differentiation gives
$W_x(j,y)=1$, $j^2+y^2=1+x^{-2}$ and
$j'^2+y'^2=1-x^{-2}+x^{-4}\leq1$ for $x\geq1$.
Cauchy--Schwarz yields $a_1,a_2$. The normalized closed kernel is
\begin{equation}
 K_B=\frac{(1+1/u)(1-1/v)-e^{-2(v-u)}(1-1/u)(1+1/v)}{2\kappa_c},
 \quad v=\kappa_cs\geq u=\kappa_cr\geq1.
\end{equation}
Both terms are nonnegative before subtraction; the first dominates,
so $0\leq K_B\leq b_0/(2\kappa_c)$. Differentiating gives the stated
$b_2$ bound (it concerns $(\partial_r-\kappa_c)K_B$, not
$\partial_rK_B$). More explicitly, twice the absolute normalized physical
derivative kernel is
\begin{equation}
 (1+u^{-1}+u^{-2})(1-v^{-1})
 +e^{-2(v-u)}(1-u^{-1}+u^{-2})(1+v^{-1})
 \leq 2+2u^{-1}+u^{-2}.
\end{equation}
The inequality follows by replacing the exponential by one, leaving
$2+2u^{-2}-2/(uv)$, which is at most the right-hand side for $u\geq1$.
The ordered-series argument now starts with a free term bounded by $b_0$.
It gives $\psi\leq b_0e^{\mu Q(r)}$ and the four errors above.
In both sectors the correction tends to zero, so
$e^{\kappa_cr}B\to1$ and the Jost solution is nonzero.
If $(E_A,E_B,E_{A'},E_{B'})$ denotes these four normalized Cauchy-data
error bounds and $s_3$ is the fixed Wronski-row scale, then
\begin{equation}
 |E_{2+i}|\leq |s_3|\bigl(E_A|A_i'(L)|+E_B|B_i'(L)|
                    +E_{A'}|A_i(L)|+E_{B'}|B_i(L)|\bigr),\quad i=1,2.
\end{equation}
The core values on the right are enclosed over the entire parameter box.
The exact dyadic scales $s_1,s_3$ change conditioning, not the zero set.

\subsection{Continuity on the parameter box}
The profile contractions use common upper bounds over the whole box,
with $\inf\kappa_0>0$ and $\sup\Lambda<1$. Their kernels vary
continuously and have a common integrable majorant. For two parameter
values $z,z'$, the fixed points obey
\begin{equation}
 \|\widetilde u_z-\widetilde u_{z'}\|_\infty
 \leq\frac{1}{1-\Lambda}
       \sup_{u\in X'}\|(T_z-T_{z'})u\|_\infty\longrightarrow0,
\end{equation}
where a common ball $X'$ contains both fixed points.
Dominated convergence also gives continuity of the profile matching
derivative. The Jost Picard iterates are continuous in parameters and
converge uniformly with the common factorial bound. Their normalized
Cauchy data are therefore continuous. Regular core solutions depend
smoothly on parameters. Consequently every component of the tail
correction $E$ in the finite-dimensional matching problem is continuous.

\subsection{Validated Taylor steps and the origin}
\label{supp:taylor}
Consider $y'=F(r,y)$ at a positive center $r_j$. Let a complex square
$\mathcal D$ of half-side $R$ enclose $|r-r_j|\leq R$, and let $B$ be a
compact product of complex rectangles. Suppose $F$ is analytic on a
neighborhood of these sets and the interval test
\begin{equation}
 Y_0+\mathcal D_0 F(r_j+\mathcal D_0,B)\subset B,
 \qquad R<r_j,
 \label{supp:taylor-test}
\end{equation}
is satisfied, where $\mathcal D_0=\mathcal D-r_j$ and $Y_0$ encloses
all initial values. If $\beta_i$ bounds the distance of component $B_i$
from its chosen center, the Taylor coefficients obey
$|y_{i,n}|\leq\beta_iR^{-n}$ for $n\geq1$. For $N\geq1$ and
$0<h<R$, $q=h/R$, the enclosed endpoint is
\begin{equation}
 y_i(r_j+h)\in\sum_{n=0}^{N-1}y_{i,n}h^n+
 \left[-\frac{\beta_iq^N}{1-q},\frac{\beta_iq^N}{1-q}\right].
 \label{supp:taylor-remainder}
\end{equation}

\noindent\emph{Proof.}
The Picard operator along straight complex segments maps analytic
functions valued in $B$ into the same set by \eqref{supp:taylor-test}.
Convexity justifies the segment average. Analyticity supplies a finite
Lipschitz constant on the enclosing set; on the disk $|r-r_j|\leq R$,
iterated differences have bound $(\Lambda_T R)^m/m!$. Picard iteration therefore converges
uniformly to the unique analytic solution. Cauchy's coefficient bound
for $y_i$ minus the center of $B_i$, first on $R'<R$ and then as
$R'\uparrow R$, gives the coefficient estimate. Summing the geometric
series gives \eqref{supp:taylor-remainder}. Interval recursion encloses
the retained coefficients. How the trial box was found does not affect
this implication: only the verified inclusion accepts the step.
The condition $R<r_j$ keeps both $2/r$ and, for $\ell=1$, $2/r^2$
analytic. It applies only to positive-center steps.

At the origin write $N(f)=(1-\omega^2)f-2f^3+\tfrac32f^5$. The
regular profile solves
\begin{equation}
 f(r)=a+r^2\int_0^1\tau(1-\tau)N(f(\tau r))\,\dd\tau,
 \qquad f'(r)=r\int_0^1\tau^2N(f(\tau r))\,\dd\tau.
\end{equation}
These formulas follow by integrating $(r^2f')'=r^2N(f)$ twice.
They give the enclosure masses $R^2/6$ and $R/3$ and the recurrence
$(n+2)(n+3)f_{n+2}=[N(f)]_n$. The derivative remainder uses its own
box for $f'$, not a differentiated endpoint error. Profile parameter
variations obey the same integral operator with source
$N'(f)\partial_pf+\partial_pN$.
For the origin step let $\mathcal D_R$ enclose the complex disk
$|r|\leq R$, and let $\mathcal D_R^{[2]}$ enclose
$\{r^2:|r|\leq R\}$. The validated a priori condition is
\begin{equation}
 a+\tfrac16\mathcal D_R^{[2]}N(B_f)\subset B_f,
 \qquad B_{f'}=\tfrac13\mathcal D_R N(B_f).
\end{equation}
It is checked componentwise over the full parameter box, with the same
condition for each differentiated source. Only this verified
self-inclusion permits the analytic Picard construction and the Cauchy
bound for the start; the derivative remainder uses $B_{f'}$.

For $\ell=0$, the mode equation is nonsingular at zero.
For $\ell=1$ set $A=r^2\alpha$, $B=r^2\beta$; the centrifugal term
cancels and each regular component satisfies $y''+4y'/r=F$ with
\begin{equation}
 y(r)=y(0)+\frac{r^2}{3}\int_0^1\tau(1-\tau^3)F(\tau r)\,\dd\tau,
 \qquad y'(r)=r\int_0^1\tau^4F(\tau r)\,\dd\tau.
\end{equation}
The enclosure masses are $R^2/10$ and $R/5$, with recurrence
$(n+2)(n+5)y_{n+2}=F_n$, $y_1=0$. For the joint source box $B$,
the corresponding verified componentwise conditions are
\begin{equation}
 y(0)+\tfrac1{10}\mathcal D_R^{[2]}F(B)\subset B_y,
 \qquad B_{y'}=\tfrac15\mathcal D_R F(B).
\end{equation}
They apply also to the parameter-variation sources. Under these
self-inclusions, Volterra iteration on the disk
has factorial convergence as above; the Cauchy remainder applies to
both $y$ and $y'$. The two regular columns have initial data
$(\alpha(0),\beta(0))=(1,0),(0,1)$ and zero derivatives.
At the first positive endpoint $h_0$, use the exact conversion
$A=h_0^2\alpha$, $A'=2h_0\alpha+h_0^2\alpha'$ (and likewise for $B$).
The parameter-variation equations use the differentiated sources and
the same origin construction. Thus the start is an enclosed regular
solution, rather than prescribed approximate data at a small radius.
For completeness, the core Jacobian calculation evolves the profile
and its $a,\omega$ derivatives, the two regular mode columns, and
their $a,\rho,\omega$ derivatives as one system of 38 real components.
In particular, with $S=f^2$,
$\partial_pS=2f\partial_pf$,
$\partial_pdV=(9S-4)\partial_pS$ and
$\partial_pC=(6S-2)\partial_pS$. The explicit frequency derivatives
are obtained from $-(\omega\pm\rho)^2$, while the profile source
has $\partial_\omega N=-2\omega f$. The parameter $c$ enters the
profile tail matching, not the core initial-value problem.

\subsection{Wronski matching, positivity and smoothness}
\label{supp:matching-positive}
For a symmetric radial matrix $M$, the bilinear form
$W(X,Y)=X\cdot Y'-X'\cdot Y$ is constant because
$W'=X\cdot MY-MX\cdot Y=0$. For $\ell=0$ choose
$R_i(0)=0$, $R_i'(0)=e_i$; then $W(\Psi,R_i)=\Psi_i(0)$.
For $\ell=1$, the Frobenius exponents are $2$ and $-1$. Choose regular
columns $R_i\sim r^2e_i$ and singular columns $S_i\sim r^{-1}e_i/3$.
Since $W(r^{-1},r^2)=3$, $W(S_i,R_j)=\delta_{ij}$ and
$W(R_i,R_j)=0$. Possible higher-order logarithmic terms do not change
these limiting pairings. Therefore $W(\Psi,R_1)=W(\Psi,R_2)=0$
removes the two singular coefficients. Profile matching equates the
inner and outer Cauchy data; ODE uniqueness then gives one global
profile and the Jost mode belongs to the regular origin span.

Suppose interval enclosures give $f>0$ on $[0,L_p]$ and
$|f|<\sqrt{S_-}$ on $[L_p,L]$, where
\begin{equation}
 S_-=(2-\sqrt{4-6\kappa_0^2})/3,
 \qquad 4-6\kappa_0^2>0.
\end{equation}
Then $V_f=\kappa_0^2-2f^2+\tfrac32f^4>0$ on that remaining interval.
With $u=rf$, $u''=V_fu$ and both endpoint values positive, a negative
interior minimum would have $u''=V_fu<0$, a contradiction. A zero
minimum has $u=u'=0$ and forces the identically zero solution, also a
contradiction. The positive exterior fixed point completes $f>0$.

The origin series for $f$ is even. For $\ell=0$, the regular mode
components $A,B$ are odd, so $A/r,B/r$ are even analytic functions.
For $\ell=1$, $\alpha,\beta$ are even analytic and
$(A/r)Y_{1m}=\alpha(r)(rY_{1m})$ is smooth, since $rY_{1m}$ is linear
in Cartesian coordinates. The ODE and the exponentially decaying
exterior bounds propagate decay to derivatives. This establishes
smooth exponentially localized spatial fields in the two stated sectors.

\subsection{Geometric radial simplicity and its extra inequality}
\label{supp:simplicity}
Assume additionally $2\kappa_0>\kappa_c$. Every solution $\Phi$ that is
square integrable on $[L,\infty)$ is a multiple of the constructed
Jost solution $\Psi$.

\noindent\emph{Proof.}
Bounded coefficients imply $\Phi''\in L^2$. The interpolation estimate
on overlapping unit intervals gives $\Phi'\in L^2$; hence
$\Phi,\Phi'\to0$. Two oscillatory comparison solutions can be
constructed by Volterra iteration about $(\cos kr,0),(\sin kr,0)$
for $\ell=0$, or $(j(kr),0),(y(kr),0)$ for $\ell=1$.
In the closed component the unnormalized Green kernel grows at most
as a constant times $e^{\kappa_c(s-L)}$, while the coupling is bounded
by a constant times $e^{-2\kappa_0s}/s^2$. The assumption makes their
product integrable. The ordered Picard series and its differentiated
series converge, with corrections tending to zero. Their mutual
Wronskian is $k\ne0$; both are Wronski-orthogonal to the decaying
$\Psi$. They and $\Psi$ span a three-dimensional subspace of the
four-dimensional space of radial Cauchy data: a nonzero combination
of the oscillatory asymptotes cannot decay. Since $\Phi,\Phi'\to0$,
$\Phi$ is Wronski-orthogonal to this whole subspace. The Wronski form
on Cauchy data is nondegenerate, so this orthogonal complement is
one-dimensional and contains $\Psi$. It is precisely its span.

The inequality $2\kappa_0>\kappa_c$ is needed for this comparison
construction, not inferred from existence of the localized mode.
For C0 the recorded lower bound on $2\kappa_0-\kappa_c$ is $0.375236$,
separately checked but omitted from the historical aggregate success
flag. The C1/C2 drivers include the condition. This appendix retains
that audit qualification and does not silently replace the C0 driver.
The conclusion is geometric simplicity in a fixed radial sector;
the $\ell=1$ angular degeneracy, algebraic multiplicity, Jordan chains
and stability are separate questions.

\subsection{Brouwer inclusion with a continuous correction}
Let $Z=z_0+[-\delta,\delta]$, $H=H_0+E$ continuous, $H_0\in C^1$,
$DH_0(Z)\subset D$ and $|E|\leq\epsilon$ componentwise. The inclusion
and weighted matrix bound in the main text imply a zero of $H$ in $Z$.
Indeed, the mean-value formula encloses $z-YH_0(z)$ by
$z_0-YH_0(z_0)+(I-YD)[-\delta,\delta]$; adding $-YE$ gives a
self-map of $Z$. Brouwer supplies $z_*$ with $YH(z_*)=0$.
Let $A=DH_0(z_0)$. The weighted row bound implies the operator norm of
$I-YA$ is less than one in the box-weighted maximum norm. Thus $YA$
is invertible, as is $Y$, and $H(z_*)=0$. Neither differentiability
of $E$ nor uniqueness of the parameter zero follows from this argument.
